% Coordinate unit: mm. Each declaration lists visible vertices, sources, then
% every directed edge, including source-to-recipient edges. Qn denotes psi_n,
% Cn denotes lambda_n, and Bp is the second output field B'. Figure panels
% use these declarations verbatim, without scaling the picture or its text.
\definecolor{censusBrown}{HTML}{886340}
\definecolor{censusAmber}{HTML}{83500D}
\definecolor{censusVisible}{HTML}{FFFBF2}
\definecolor{censusGreen}{HTML}{245E43}
\definecolor{censusGreenLine}{HTML}{42805A}
\definecolor{censusMint}{HTML}{E5F3E7}
\definecolor{censusRed}{HTML}{913E46}
\definecolor{censusRedLine}{HTML}{AF5960}
\definecolor{censusRose}{HTML}{FBE8E7}
\definecolor{censusInk}{HTML}{41454A}
\tikzset{
  census node/.style={inner sep=0pt,outer sep=.23pt,minimum size=8mm,line width=.55pt,
    font=\normalfont\fontsize{10.5}{11}\selectfont},
  census visible/.style={census node,rectangle,rounded corners=.8mm,
    draw=censusBrown,
    fill=censusVisible,text=censusAmber},
  census latent/.style={census node,circle},
  census Q/.style={census latent,draw=censusGreenLine,fill=censusMint,text=censusGreen},
  census C/.style={census latent,draw=censusRedLine,fill=censusRose,text=censusRed},
  census edge/.style={-{Latex[length=1.65mm,width=1.08mm]},line width=.58pt,
    shorten >=.25pt,shorten <=.3pt}
}
\newcommand{\censusLabelBp}{\mathbf B'}
\newcommand{\censusVisibleLabel}[1]{%
  \ifcsname censusLabel#1\endcsname\csname censusLabel#1\endcsname
  \else\mathbf{#1}\fi}
\def\censusSymbolQ{\psi}
\def\censusSymbolC{\lambda}
\def\censusLineQ{censusGreenLine}
\def\censusLineC{censusRedLine}
\def\censusSource#1#2\relax#3#4{%
  \node[census #1] (#1#2) at (#3,#4) {$\csname censusSymbol#1\endcsname_{#2}$};}
\newcommand{\censusPlot}[3]{%
  \begin{tikzpicture}[x=1mm,y=1mm,baseline=(current bounding box.center)]
    \path[use as bounding box] (-3,-4) rectangle (64,48);
    \foreach \v/\xx/\yy in {#1}{
      \node[census visible] (\v) at (\xx,\yy) {$\censusVisibleLabel{\v}$};}
    \foreach \v/\xx/\yy in {#2}{\expandafter\censusSource\v\relax{\xx}{\yy}}
    \begin{scope}[on background layer]
      \foreach \tail/\head in {#3}{
        \def\edgecolor{censusInk}%
        \expandafter\censusEdgeColor\tail\relax
        \draw[census edge,draw=\edgecolor] (\tail) -- (\head);}
    \end{scope}
  \end{tikzpicture}}
% Visible C and classical source Cn share an initial. Only sources have a suffix
% after Q or C. The primed visible Bp always uses the visible-arrow color.
\def\censusEdgeColor#1#2\relax{%
  \if\relax\detokenize{#2}\relax\else
    \ifcsname censusLine#1\endcsname\edef\edgecolor{\csname censusLine#1\endcsname}\fi
  \fi}
\newcommand{\Diagram}[4]{%
  \expandafter\newcommand\csname censusPanel#1\endcsname{\censusPlot{#2}{#3}{#4}}}
\newcommand{\censusPanel}[2]{%
  \begin{minipage}[c]{.46\linewidth}\centering
    \csname censusPanel#1\endcsname\par\vspace{.7mm}
    {\small\sffamily #2\par}
  \end{minipage}}
\newcommand{\censusPair}[6]{%
  \expandafter\newcommand\csname censusFigure#1\endcsname{%
    \begin{figure}[htbp]
      \centering
      \censusPanel{#2}{\textbf{(a)}\enspace #3}\hfill
      \makebox[0pt]{\raisebox{2ex}{$\longrightarrow$}}\hfill
      \censusPanel{#4}{\textbf{(b)}\enspace #5}
      \caption{#6}\label{fig:#1}
    \end{figure}}}
\newcommand{\censusfigure}[1]{\csname censusFigure#1\endcsname}

\Diagram{deg-before}
  {C/30/38,B/8/7,S/52/7}
  {Q1/9/29,Q3/51/29,Q2/30/5}
  {Q1/B,Q1/C,Q3/C,Q3/S,Q2/B,Q2/S}
\Diagram{deg-after}
  {C/30/38,B/8/7,S/52/7}
  {C1/9/29,C3/51/29,Q2/30/5}
  {C1/B,C1/C,C3/C,C3/S,Q2/B,Q2/S}
\censusPair{degradation}{deg-before}{All-quantum triangle}
  {deg-after}{Two sources made classical}
  {Degrading \(\psi_1\) on \(BC\) and \(\psi_3\) on \(CS\) gives the
   classical sources \(\lambda_1,\lambda_3\). Laboratory \(C\) records
   both settings, while \(B,S\) measure \(\psi_2\). This is the final
   step after the upward lookup in Figure~\ref{fig:upward} and fixing \(A\).}

\Diagram{fix-before}
  {A/10/36,B/13/8,C/35/8,D/54/30}
  {Q1/2/21,Q2/32/31}
  {A/B,B/C,C/D,Q1/A,Q1/B,Q2/A,Q2/C,Q2/D}
\Diagram{fix-after}
  {B/13/8,C/35/8,D/54/30}
  {Q2/32/31}
  {B/C,C/D,Q2/C,Q2/D}
\censusPair{fixing}{fix-before}{Two source inputs at \(A\)}{fix-after}{Fix \(A\)}
  {Fixing \(A\) to a constant removes its message to \(B\).
   Source \(\psi_2\) now supplies only \(C,D\), and the surviving share of
   \(\psi_1\) can be prepared locally at \(B\). The setting \(B\), wire
   \(C\to D\), and shared state on \(CD\) give an instrumental seed.}

\Diagram{naive-before}
  {A/8/8,B/52/8,C/30/26,D/30/43}
  {Q1/30/4,Q2/8/28,Q3/52/28}
  {C/D,Q1/A,Q1/B,Q2/A,Q2/C,Q3/B,Q3/C,Q3/D}
\Diagram{naive-after}
  {A/8/8,B/52/8,D/30/43}
  {Q1/30/4,C2/8/28,Q3/52/28}
  {Q1/A,Q1/B,C2/A,C2/D,Q3/B,Q3/D}
\censusPair{naive}{naive-before}{Classical message \(C\to D\)}
  {naive-after}{Marginalize \(C\)}
  {Ordinary marginalization makes a classical register from the \(AC\)
   source available to \(A,D\), drawn as \(\lambda_2\). The original
   \(BCD\) source retains the \(B,D\) shares and can also carry its own
   forwarded classical register. Degrading this last quantum source gives
   the triangle seed with setting recorder \(D\).}

\Diagram{tele-before}
  {A/8/35,B/29/22,C/9/6,D/52/9}
  {Q1/48/35,Q2/2/21}
  {A/B,B/C,B/D,Q1/A,Q1/B,Q1/D,Q2/A,Q2/C}
\Diagram{tele-after}
  {B/29/22,C/9/6,D/52/9}
  {Q1/48/35,Q2/2/21}
  {B/C,B/D,Q1/B,Q1/D,Q2/B,Q2/C}
\censusPair{teleportation}{tele-before}{Entanglement shared by \(A,B\)}
  {tele-after}{Marginalize \(A\)}
  {The source \(\psi_1\) on \(ABD\) can supply an entangled pair to
   \(A,B\) alongside the desired state for \(B,D\). Laboratory \(A\)
   uses the pair to teleport its share of \(\psi_2\) to \(B\) and
   sends the correction along \(A\to B\). The two states on \(BC\)
   and \(BD\) give the Evans seed.}

\Diagram{stitch-before}
  {A/52/38,B/30/28,C/8/15,D/52/15}
  {Q1/16/3,Q2/44/3}
  {A/D,B/C,Q1/B,Q1/C,Q2/B,Q2/D}
\Diagram{stitch-after}
  {B/30/28,C/8/15,D/52/15}
  {Q1/16/3,Q2/44/3}
  {B/C,B/D,Q1/B,Q1/C,Q2/B,Q2/D}
\censusPair{stitching}{stitch-before}{Outcome \(B\) and setting \(A\)}
  {stitch-after}{Stitch \(A\) onto \(B\)}
  {The stitch retains the wire \(B\to C\) and redirects \(A\to D\)
   to \(B\to D\). Keeping records with \(A=B\) supplies the Evans
   experiment on the right, with two independent sources incident to \(B\).}

\Diagram{post-before}
  {B/7/38,C/14/9,D/49/9,A/46/38}
  {Q1/25/26,Q2/57/25}
  {A/B,B/C,C/D,Q1/A,Q1/C,Q2/A,Q2/D}
\Diagram{post-after}
  {B/7/38,C/14/9,D/49/9}
  {Q3/33/28}
  {B/C,C/D,Q3/C,Q3/D}
\censusPair{postselection}{post-before}{Two independent states at \(A\)}
  {post-after}{Postselect at \(A\)}
  {Postselecting a joint-measurement outcome at \(A\) leaves \(C,D\)
   with a shared state, drawn as \(\psi_3\). On the selected branch
   \(B\) receives a fixed value from \(A\). It supplies a setting to \(C\)
   in the resulting instrumental experiment.}

\Diagram{split-before}
  {A/12/23,B/26/42,C/36/23,D/47/23}
  {Q1/0/23,Q2/24/23,Q3/60/23}
  {A/B,B/C,C/D,Q1/A,Q1/B,Q2/A,Q2/C,Q3/B,Q3/D}
\Diagram{split-after}
  {A/12/23,B/26/42,Bp/26/3,C/36/23,D/47/23}
  {Q1/0/23,Q2/24/23,Q3/60/23}
  {A/B,A/Bp,B/C,Bp/C,C/D,Q1/A,Q1/B,Q1/Bp,Q2/A,Q2/C,Q3/B,Q3/Bp,Q3/D}
\censusPair{splitting}{split-before}{One reported field at \(B\)}
  {split-after}{Split \(B\) into two fields}
  {Splitting \(B\) into \(B,B'\) gives both fields the original inputs
   of \(B\) and sends both outputs to \(C\). Either source enlarged by the
   split can supply the shared private randomness. Panel (b) shows the
   intermediate graph before parent prediction restricts \(B'\).}

\Diagram{parent-before}
  {A/12/23,B/26/42,Bp/26/3,C/36/23,D/47/23}
  {Q1/0/23,Q2/24/23,Q3/60/23}
  {A/B,A/Bp,B/C,Bp/C,C/D,Q1/A,Q1/B,Q2/A,Q2/C,Q3/B,Q3/D}
\Diagram{parent-after}
  {B/16/30,C/35/15,D/57/27,Bp/14/0}
  {Q2/35/33,Q3/38/44,C2/14/15}
  {B/C,Bp/C,C/D,Q2/B,Q2/C,Q3/B,Q3/D,C2/B,C2/Bp,C2/C}
\censusPair{parent}{parent-before}{Only \(A\) supplies \(B'\)}
  {parent-after}{Marginalize predictor \(A\)}
  {Laboratory \(A\) predicts \(B'\). Panel (a) shows the intermediate
   graph obtained by removing the two source inputs to \(B'\) from
   Figure~\ref{fig:splitting}(b). Marginalizing \(A\) gives panel (b).
   Laboratory \(A\) uses \(\psi_1\) to teleport the share of \(\psi_2\) to
   \(B\) and forwards a classical register from the \(AC\) source, giving
   \(\lambda_2\) on \(B,B',C\). The next census move marginalizes \(B\)
   to reach an instrumental seed.}

\Diagram{dsep-before}
  {A/8/41,B/40/31,C/52/3,D/20/13}
  {Q1/52/41,Q2/20/31,Q3/8/3,Q4/40/13}
  {Q1/A,Q1/B,Q1/C,Q2/A,Q2/B,Q2/D,Q3/A,Q3/C,Q3/D,Q4/B,Q4/C,Q4/D}
\Diagram{dsep-after}
  {B/30/38,C/8/7,D/52/7}
  {C1/9/29,C2/51/29,Q3/30/5}
  {C1/B,C1/C,C2/B,C2/D,Q3/C,Q3/D}
\censusPair{dsep}{dsep-before}{Four independent triple sources}
  {dsep-after}{Restrict \(B\) and forget \(A\)}
  {Laboratory \(A\) predicts \(B\) using the sources \(\psi_1,\psi_2\).
   The d-separation test allows \(B\) to drop the input from \(\psi_4\)
   and to receive only classical data from \(\psi_1,\psi_2\). Removing
   \(A\) leaves both \(\psi_3\) and \(\psi_4\) on \(CD\). Panel (b) draws
   the tensor product as \(\psi_3\).}

\Diagram{kpc-before}
  {C/9/37,D/9/9,E/50/37,F/50/9}
  {Q1/57/23,Q2/30/5}
  {C/D,C/E,D/E,Q1/E,Q1/F,Q2/D,Q2/F}
\Diagram{kpc-after}
  {C/9/37,D/9/9,E/50/37,F/50/9}
  {C1/57/23,Q2/30/5}
  {C/D,C1/E,C1/F,Q2/D,Q2/F}
\censusPair{semigraphoid}{kpc-before}{\(F\) predicts \(E\)}
  {kpc-after}{Restrict the inputs of \(E\)}
  {When \(F\) predicts \(E\), the semigraphoid certificate allows
   deletion of \(C\to E\) and \(D\to E\). In the Bell graph, \(\lambda_1\)
   supplies the setting at \(E,F\), and \(\psi_2\) supplies the state
   measured at \(D,F\). The proof constructs a fresh source on \(EF\)
   in the classical model.}
\censusPair{entropic}{kpc-before}{Prediction constraint \(H(E\mid F)=0\)}
  {kpc-after}{Bell experiment}
  {The entropy program certifies the deletion of \(C\to E\) and
   \(D\to E\), the same census move as Figure~\ref{fig:semigraphoid}.
   The proof uses \(H(E\mid F)=0\) and the observable independence
   \(I(E;CD)=0\) to justify a fresh classical source on \(EF\). The
   census records the semigraphoid certificate for the move.}

\Diagram{up-before}
  {A/30/28,B/30/7,C/9/38,S/51/38}
  {Q1/12/24,Q2/48/24,Q3/30/42}
  {A/B,Q1/A,Q1/B,Q1/C,Q2/A,Q2/B,Q2/S,Q3/C,Q3/S}
\Diagram{up-after}
  {A/30/28,B/30/7,C/9/38,S/51/38}
  {Q1/12/24,Q2/48/24,Q3/30/42}
  {Q1/A,Q1/B,Q1/C,Q2/A,Q2/B,Q2/S,Q3/C,Q3/S}
\censusPair{upward}{up-before}{Target with \(A\to B\)}
  {up-after}{Remove the visible arrow}
  {The two graphs share every source and have the same classical model set.
   The quantum model set on the right is strictly smaller. An upward
   lookup transfers a gap from right to left. Fixing \(A\) on either side
   gives an all-quantum triangle on \(B,C,S\).}
